% !TeX root = ../main.tex
\section{Background and Related Work}

This section is a scoped narrative synthesis selected to establish context for ArchSync's design, not a systematic literature review. Sources were drawn from the retained keyword-search corpus and selected for their relevance to architecture conformance, drift detection, and static-analysis reliability, with emphasis on 2022--2025 publications. Four earlier works are retained as concept-origin references. Selection was purposive rather than exhaustive; the discussion distinguishes evaluated techniques, tool demonstrations, and secondary syntheses, and does not establish the absence of competing approaches.

\subsection{Architecture Erosion and Reconstruction}

Architecture erosion denotes progressive misalignment between intended architecture and implementation. A mapping study of 73 primary studies characterizes its manifestations and consequences \cite{li2022erosion}. Evidence from two OpenStack projects identifies erosion symptoms in code reviews \cite{li2022symptoms}; a subsequent journal study examines violation symptoms and developer responses across OpenStack and Qt projects \cite{li2023warnings}. These related studies offer concrete developer-facing evidence, but their overlapping project contexts should not be counted as independent replications. A mixed-methods study further identifies documentation, clear responsibilities, and shared practices as ways developers manage drift \cite{anthony2024drifting}. Together, these findings motivate implementation-linked checks without reducing architecture quality to one structural metric.

Reconstruction recovers architectural views from implementation artifacts. Ducasse and Pollet's foundational taxonomy distinguishes bottom-up, top-down, and hybrid processes \cite{ducasse2009reconstruction}. A recent tertiary mapping organizes the applicability of recovery approaches, emphasizing that selecting an appropriate method depends on its context and assumptions \cite{sar2024taxonomy}. A microservice-focused mapping compares the purposes, inputs, and outputs of reconstruction tools and identifies availability and maintenance concerns \cite{bakhtin2024mapping}. These syntheses provide context for method selection rather than a uniform benchmark of tool accuracy. View-based drift analysis demonstrates how reconstructed specifications can be compared with intended models in an industrial case study \cite{uzun2024drift}. ArchSync adopts a narrower observation boundary: an explicit expected model and a defined set of service and infrastructure relationships.

\subsection{Static Architecture Conformance}

Software Reflexion Models compare a high-level model with a source-derived model and distinguish convergence, divergence, and absence \cite{murphy1995reflexion}. Knodel and Popescu compare several static conformance formulations \cite{knodel2007comparison}, while Dependency Constraint Language expresses allowed and forbidden dependencies in object-oriented systems \cite{terra2009dcl}. These works supply the conceptual origins of expected-versus-observed comparison and dependency rules; neither concept is claimed as new here.

Recent tools adapt conformance checking to particular languages and deployment settings. CALint checks Python implementations against the Clean Architecture dependency rule and accompanies detection with refactoring support \cite{calint2022}. ArchLintor is a particularly close language-level comparison: it detects architectural violations in TypeScript systems, reports module relationships, and supports CI/CD integration \cite{archlintor2025}. ArchSync focuses on typed service-topology observations and a versioned finding contract, whereas this ArchLintor description centers on module-interaction checking. No head-to-head comparison with ArchLintor is reported.

IaC coupling-conformance metrics extend architecture checking to deployment practices \cite{ntentos2022iac}, providing a relevant reference for ArchSync's planned, unevaluated IaC extension. In a correction-oriented approach, 3Erefactor searches for executable refactorings that reduce architectural inconsistency \cite{refactor2024}; ArchSync currently reports findings and gate outcomes rather than generating repairs. Language Models as Architectural Gatekeepers explores converting natural-language design rules into executable conformance tests \cite{lucena2025gatekeepers}. Its preliminary evaluation distinguishes tests that execute from tests that correctly express the intended rule. ArchSync instead evaluates explicit rules deterministically. This design avoids dependence on generated rules at decision time, but determinism alone does not establish detector accuracy or superiority over an LLM-assisted approach.

The prototype combines typed HTTP, data, and asynchronous edges; explicit require and require-path semantics; source and model evidence in a versioned finding contract; and separate violation and evolution outcomes. This is an implementation-specific positioning within the reviewed work, not an exhaustive novelty claim.

\subsection{Change-Time Conformance}

Architectural change can be examined at several time scales. Commit-level prediction relates source changes to architectural changes \cite{commits2025}; a predicted change is not necessarily a violation of intended architecture. Repository mining combined with static reconstruction tracks microservice architectural evolution across development history \cite{evolution2024}. View-based drift analysis compares recovered and intended specifications \cite{uzun2024drift}, while an industrial conformance-visualization study reports practitioner feedback and challenges involving scalability, usability, and workflow integration \cite{thermofisher2025}. These approaches motivate change-aware feedback but evaluate different questions and should not be treated as interchangeable baselines.

ArchSync places its comparison at pull-request time and exposes deterministic findings through annotations and exit-code decisions. Its current evidence verifies change-scoped behavior on a co-developed benchmark. It does not establish industrial adoption, reduced developer effort, organizational effectiveness, or comparative performance against the reviewed tools.

\subsection{Evaluation and Static-Analysis Reliability}

Konersmann et al. examine evaluation and replication support in software-architecture research \cite{konersmann2022replicability}. Abgaz et al.'s review concerns monolith-to-microservice decomposition \cite{abgaz2023decomposition}; that adjacent scope should not be treated as a direct measurement of architecture-conformance accuracy or proof that suitable conformance baselines do not exist. A more direct precedent is Schneider et al.'s comparison of static microservice-recovery tools: it identifies 13 tools, executes nine, and reports both individual and combined-tool recovery performance \cite{schneider2025comparison}. This provides a concrete model for future external-baseline work. ArchSync has not reproduced that benchmark, and its typed-edge detector results cannot be compared numerically without aligning capabilities and ground truth.

Static-analysis-based visual models expose microservice dependencies at a system level \cite{staticmodels2024}. Recovery from multiple dependency models explores how complementary relations affect reconstructed structure \cite{dependencies2024}. These approaches motivate the future multi-source direction discussed for ArchSync, while leaving its current evidence restricted to the declared TypeScript source patterns. Secondary mappings and tertiary taxonomies help organize this space; they are not additional independent trials of each underlying tool.

False positives and false negatives remain central to reliability. Interrogation testing distinguishes precision and soundness issues and demonstrates failures in mature analyzers \cite{kaindlstorfer2024interrogation}. ArchSync consequently includes annotated positive and hard-negative signals alongside topology changes. Its manifests, pinned artifacts, explicit denominators, and replay checks support reproducibility, but do not turn the synthetic development benchmark into representative evidence. The narrative synthesis establishes design context and comparison dimensions; it supports neither coverage-completeness claims nor absence-of-evidence claims.
